% Lösungen Einheit 1 von 5 – Fläche zwischen Graph und x-Achse (zusammenbau v0.1)
\einheitenkopf{Einheit 1 von 5 · Fläche zwischen Graph und x-Achse}

\erg{9}{a) ganz oberhalb der x-Achse – ein Integral genügt \quad b) Nullstellen $0$ und $4$; Integralwert $-\frac{32}{3}$; $A = \frac{32}{3} \approx 10{,}67$ \quad c) Nullstellen $1$, $2$, $3$; $A = \frac{1}{4} = 0{,}25$ \quad d) Nullstelle rechts $4$; $A = \frac{80}{3} \approx 26{,}67$ \quad e) Nullstellen $-3$, $0$, $3$; Teilstück von $0$ bis $3$: $\left|-\frac{81}{4}\right|$; wegen der Punktsymmetrie $A = 2 \cdot \frac{81}{4} = 40{,}5$ \quad f) $\frac{128}{3}$ Flächeneinheiten, eine Flächeneinheit $= 0{,}25\,\text{m}^2$; Wand $= \frac{32}{3} - 2 \approx 8{,}67\,\text{m}^2$ \quad g) Nein: Das Dreieck mit den Ecken $(0|0)$, $(6|0)$ und dem Tiefpunkt $(3|-4{,}5)$ liegt ganz in der Fläche, sein Inhalt ist $\frac{1}{2} \cdot 6 \cdot 4{,}5 = 13{,}5 > 12$.}
\erg{10}{a) Fehler: Der negative Integralwert ist keine Fläche. Richtig: $A = \left|-\frac{4}{3}\right| = \frac{4}{3}$ \quad b) Unterhalb der x-Achse sind die Funktionswerte negativ, das Integral zählt die Fläche also negativ – ein Flächeninhalt ist aber nie negativ. \quad c) $A = \frac{85}{6} \approx 14{,}17$ Flächeneinheiten, real $\approx 1\,416{,}7\,\text{m}^2$; Säcke $= 1\,416{,}7 : 250 \approx 5{,}7$, also $6$ Säcke}
